% !PRINTED 31
\weberpage{31}
%<<CONT
is set, the differential (5) passes into
\weq{8}{\frac{1}{M}\,\frac{\xi\,d\eta - \eta\,d\xi}{\sqrt{\varphi(\xi,\eta)}}}
It can now be shown that, as soon as condition (6) is satisfied, $H$ is
divisible by $T$, and hence, since the degree of the two functions is the
same, $M$ becomes a constant. This constant is called
the \gesperrt{multiplier} of the transformation.

For if we note that, since $f(x,y)$ contains no square factor, the two
functions
\[
\frac{\partial f}{\partial x}, \quad \frac{\partial f}{\partial y}
\]
have no common divisor, and that in consequence, because $U$ and $V$ are
coprime, the functions
\[
\frac{\partial f}{\partial U}, \quad \frac{\partial f}{\partial V}
\]
likewise have no common divisor (in relation to $\xi,\eta$), then two
functions $\alpha,\beta$ of $\xi,\eta$ can be so determined that
\[
\alpha\,\frac{\partial f}{\partial U} + \beta\,\frac{\partial f}{\partial V}
\]
is coprime to an arbitrarily given function $T$. One may, e.\,g., take
$\alpha$ coprime to $T$ and $\beta$ divisible by those divisors of $T$
which do not divide $\partial f/\partial U$, but indivisible by the
common divisors of $T$ and $\partial f/\partial U$. If we now
differentiate equation (6), which we assume to be satisfied, with respect
to $\xi,\eta$, there follows:
\[
\begin{aligned}
\frac{\partial f}{\partial U}\frac{\partial U}{\partial \xi} + \frac{\partial f}{\partial V}\frac{\partial V}{\partial \xi} &= TX,\\
\frac{\partial f}{\partial U}\frac{\partial U}{\partial \eta} + \frac{\partial f}{\partial V}\frac{\partial V}{\partial \eta} &= TY,
\end{aligned}
\]
and from this, by solution with respect to $\partial f/\partial U$,
$\partial f/\partial V$:
\[
H\left(\alpha\,\frac{\partial f}{\partial U} + \beta\,\frac{\partial f}{\partial V}\right) = TZ,
\]
in which $X,Y,Z$ are integral homogeneous functions of $\xi,\eta$. But
from the last equation one concludes that $H$ must be divisible by $T$.
